\documentclass[11pt]{article}
\usepackage[a4paper,margin=27mm]{geometry}
\usepackage{fontspec,amsmath,amssymb,amsthm,booktabs,tabularx,microtype}
\usepackage[round,authoryear]{natbib}
\usepackage[hidelinks]{hyperref}
\usepackage{needspace}
\AddToHook{cmd/section/before}{\Needspace{10\baselineskip}}
\widowpenalty=10000
\clubpenalty=10000
\brokenpenalty=10000
\setmainfont{TeX Gyre Pagella}
\setmonofont{Latin Modern Mono}[Scale=0.86]
\setlength{\emergencystretch}{2em}
\newtheorem{proposition}{Proposition}
\newcommand{\Accept}{\mathsf{Accept}}
\newcommand{\Valid}{\mathsf{Valid}}
\title{Theorem Notary: A Glass-Box Protocol for Publishing and Reusing Kernel-Checked Theorems}
\author{Alex Chengyu Li\\\small Independent Researcher}
\date{}
\begin{document}
\maketitle
\begin{abstract}
Formal proof libraries support theorem reuse, while independently maintained
projects must also preserve which exact statement, proof data, and logical
assumptions they have selected. We propose Theorem Notary, a mathematical
engineering protocol for this publication boundary. Its Glass Box interface
makes a theorem usable as a component while keeping its proof, definitions,
and dependencies available for inspection and checking by the kernel, the proof assistant's logical core. The architecture
separates provider assembly, immutable publication, consumer-controlled
proof admission, and local acceptance reuse. A certificate binds one declaration to
fixed source and proof data; a consumer lock fixes the selected identities.
Provider signatures attribute supply. Admission validates the actual proof
objects and every inherited certificate under a disclosed foundation.
A concrete Lean profile specifies the encoding, dependency bindings, data-only
checking process, authenticated local records, and build-failure behavior.
We explain the assurance supplied by each stage and give an abstract composition
argument. Evaluation combines a large archived formalization containing over
30,000 modules and about 121 GiB of unpacked proof data with a fresh three-project
case in which later proofs consume separately certified conclusions.
Adversarial evaluation attempts to substitute statements under valid signatures,
conceal inner proof obligations, and forge consumer acceptance. Each attempt
is tied to the check that detects or prevents it. The protocol gives successive consumers an executable contract for
preserving exact theorem identity, supply attribution, and checking scope.
This supports the accumulation of reusable formal components as AI-assisted
research increases demand on formalization and review.
\end{abstract}

\section{From understanding a theorem to handing it on}

A mathematician considering an external theorem has good reasons to begin
with its statement and proof. The quantifiers must match the intended problem,
the definitions must express the intended objects, and the axioms must be
acceptable for the intended use. Reading the argument also supplies ideas that
can guide a new proof. These judgments remain central when a formal proof is
available. A publication protocol should preserve the researcher's ability to
make them and make their consequences reliable in subsequent work.

Ordinary formal imports already provide the logical step of reuse. A theorem
proved in an external Lean project can participate in a downstream proof;
Lean's replay utility can additionally validate an environment by sending its
declarations back to the kernel~\citep{leanreplay}. A careful maintainer can
pin dependencies, inspect actual axiom dependencies, and maintain a controlled
build. The case for a further protocol begins when an accepted result must
travel between projects whose maintainers make these decisions separately.

Consider a first researcher who publishes a theorem, a second who imports it
and proves two new lemmas, and a third who needs one of those lemmas. At the
first handoff, the collaborators may agree directly on a source version,
statement, and checking basis. At the next handoff, the third researcher needs
to recover what was actually selected and which parts of that basis remain
applicable. If an inner component changes, someone must decide whether to
retain the earlier source and proof data or adopt the replacement. A successful build
answers a logical question within its current environment; the research
collaboration also needs an explicit account of what that environment is
supposed to preserve.

The distinction matters even when every proof is valid. A statement about all
positive integers and one about sufficiently large integers can both have
correct proofs under the same axioms. A printed expression can also retain a
familiar predicate name while the definition behind that name changes. With
an unchanged, fully specified downstream goal, the kernel rejects an invalid
inference from a replacement. When definitions, goals, or publication choices
change together, kernel checking establishes correctness of the resulting
formal development. Determining whether it is still the result previously
chosen requires comparing its identity and interpretation. The axiom list
answers which logical assumptions a proof uses; it leaves this continuity
question to the surrounding workflow.

Human review and identity checking therefore have complementary jobs. The
researcher decides that a precise statement is relevant and that its foundation
is suitable. The protocol records which formal realization, the concrete source and proof data, was selected.
It makes replacement visible and enforces the consumer's declared conditions
for accepting the component during use. A later change can be evaluated as a new choice. This makes a careful
judgment actionable after the original conversation has ended. Interpreting
the theorem remains a mathematical task, including checking its correspondence
to the informal problem.

Direct import can be accompanied by all of these measures. A team can maintain
exact dependency locks, proof and source inventories, declaration-level
assumption policies, authenticated attribution, and rules for reusing earlier
checks. Such a workflow already performs the work addressed here. A common
protocol gives those measures a shared meaning at the next handoff: which
object was selected, what a check covered, and which changes require a fresh
acceptance check. In a stable project under one maintainer, an established import
workflow may be sufficient. The protocol addresses repeated publication and
reuse across independently maintained projects, where otherwise each boundary
must reconstruct these agreements.

AI-assisted research gives this coordination problem particular urgency.
Producing a candidate statement or argument leaves the work of constructing a
complete formal proof, checking its intended meaning, and deciding how it can
be used. As candidate production expands, these tasks can become the limiting
capacity. An effective workflow must concentrate that capacity on the new
argument while making completed formalizations available for further work.
Ordinary imports already make such reuse possible within compatible projects.
The additional engineering task is to preserve an explicit, inspectable basis
for that reuse as components pass between maintainers and automated consumers.
A protocol serves this task by making the conditions of reuse executable.

Theorem Notary expresses this contract through per-declaration certificates.
Each certificate binds a theorem's full interface, its proof realization, and
its dependencies. A consumer lock fixes the exact identities that the consumer
has chosen. Admission then checks the supplied proof data and its inherited
certificates against those choices. Publicly inspectable proof material,
definitions, and checking conditions give the contract its Glass Box character.
The chain can grow while each consumer retains an explicit basis for deciding
what to accept.

Two records remain distinct throughout. Provider credit identifies the key
that supplied a component. Mathematical acceptance records the consumer's
check of the actual formal object under a disclosed foundation. Anyone can
supply a component; acquiring a signature never creates a new logical premise.
A fresh consumer performs its own admission, and an existing consumer can reuse
its own authenticated record for an unchanged checking scope. What travels
between researchers is a fully specified publication and its inspectable proof
basis. The decision to accept it remains with the receiving project.

This separation supports a useful form of mathematical specialization. One
researcher can provide a substantial checked component; another can select an
auxiliary lemma, build on it, and publish the resulting theorem with the
inherited identities intact. A registry can make these publications discoverable
without deciding their mathematical admissibility. Ordinary external imports
supply the logical reuse mechanism; the protocol supplies a consistent contract
for that mechanism across publication boundaries. The distinction between
closing a proof and making a closed result durably composable motivates the
mathematical engineering layer~\citep{engineering}. Here it becomes a concrete
research question: how can independently published theorems retain their exact
identity, credit, and checking basis through successive derivations?

\section{Protocol objects and participants}

The participants are a provider, a publication registry or mirror, and a
consumer. A provider assembles an already checkable component and publishes a
credit record. A registry supplies discovery and distribution. A consumer
chooses the component it intends to use and applies its own acceptance policy.
One person may play several roles, and several providers may supply the same
component. Registry membership and provider reputation are not premises of
the mathematical acceptance rule.

Glass Box describes the relationship between an interface and its
justification. Like a software package, a theorem component offers an interface
that a client can use without reproducing the implementation. Here the
interface states a proposition, and the implementation includes its proof.
The client can reason with the exported theorem while its justification remains
available for inspection: the proof objects, the meanings of referenced
constants, all recursively referenced dependencies, and the declared logical foundation. The
kernel checks that justification. A provider's claim that the interface works
supplies no proof.

The admissible components are \emph{Kernel-Only}: exported statements have
complete proof objects checkable by the named kernel under an explicitly
permitted foundation. Every additional theorem obligation must be discharged.
An arbitrary assertion cannot acquire axiom status through signing or registry
inclusion. The foundational axioms themselves are disclosed and remain subject
to the consumer's policy. Ordinary Lean imports already support this form of
proof-based abstraction. The protocol preserves its inspectable justification
across separately published components and their derived publications.

The protocol uses the objects in Table~\ref{tab:objects}. A \emph{proof
realization} is a fixed collection of source and proof data under a named
checking profile, which specifies how the component is to be checked.
A \emph{declaration certificate} selects an actual theorem
within that realization and records its full interface. The interface includes
the declaration name, type expression, universe parameters describing levels
in Lean's hierarchy of types, and the identity of the environment that gives its referenced constants meaning. Interface comparison uses these structural fields even when display formulas
coincide.

\begin{table}[ht]
\centering\small
\begin{tabularx}{\linewidth}{@{}lX@{}}
\toprule
Object & Protocol function\\
\midrule
Proof realization & Fixes source, proof bytes, dependency publications, and checking profile.\\
Declaration certificate & Binds one exact theorem interface and allowed foundation to a realization.\\
Provider credit & Signs the certificate identity and records a provider key and declared name.\\
Consumer lock & Fixes the component identities selected independently by the consumer.\\
Acceptance record & Records the consumer's actual check and the scope within which it may be reused.\\
\bottomrule
\end{tabularx}
\caption{Distinct objects for publication, selection, and mathematical admission.}
\label{tab:objects}
\end{table}

A realization may support many declaration certificates. An auxiliary lemma
and a repository's headline theorem are equally eligible if their actual
proofs satisfy the profile. Their certificates remain distinct even when they
share every module byte. A downstream author can therefore select several
conclusions from one publication and later publish new conclusions that use
them. The identity of the publication container does not determine the
granularity of mathematical reuse.

Object identities are content digests computed from a unique byte representation
with a different prefix for each object kind. A parent realization binds both the realization identities and the
certificate identities of its dependency publications. Binding only a package
identity would leave the selection of its exported interfaces underspecified.
Credit lists are outside mathematical identity, so adding another valid
provider signature preserves the selected theorem. Mutable discovery metadata
may point to these immutable identities, but it cannot update a consumer lock.

A checking profile fixes the proof language, toolchain, permitted foundation,
serialization conventions, and admission procedure. The protocol does not
infer mathematical equivalence from changes in notation or software version.
A changed realization obtains a new identity. Transport between profiles
requires a separately justified mechanism; a provider signature does not
perform that transport.

\section{Technical architecture and acceptance contract}
\label{sec:architecture}

\subsection{Responsibilities and control boundaries}

Consider the second researcher in the handoff of Section~1. The incoming
publication is controlled by its provider, while the receiving project owns
the choice of theorem and the conditions for using it. The architecture makes
that division executable. Table~\ref{tab:architecture} assigns an input, output,
and acceptance obligation to each stage. A \emph{publication bundle} is the
transport object containing a realization, its declaration certificates and
credits, and the inherited publications. The receiver can inspect those public
objects without a provider account or a private registry response.

\begin{table}[ht]
\centering\small
\renewcommand{\arraystretch}{1.12}
\begin{tabularx}{\linewidth}{@{}>{\raggedright\arraybackslash}p{0.22\linewidth}>{\raggedright\arraybackslash}X>{\raggedright\arraybackslash}X@{}}
\toprule
Stage and controller & Input and output & Required obligation\\
\midrule
Assembly: provider & Source and selected declarations produce a proof realization. & Build and check the actual exported theorems and inherited components.\\[3pt]
Publication: provider; distribution: mirror & Realization and certificates form a portable bundle with credits. & Bind exact bytes, interfaces and dependency publications; authenticate supply.\\[3pt]
Selection: consumer & A reviewed publication supplies the identities entered in a consumer lock. & Keep selection independent of the incoming object's authority.\\[3pt]
Admission: consumer & Bundle, lock and checking policy produce rejection or mathematical acceptance. & Compare every certificate with the actual checked declaration and permitted foundation.\\[3pt]
Use and republication: consumer becoming provider & Accepted components enter a derived proof and a new publication. & Preserve checked inputs through use and carry inherited identities into the new bundle.\\
\bottomrule
\end{tabularx}
\caption{The technical handoffs and the party responsible for each. Registry
availability affects delivery; registry approval supplies no proof premise.}
\label{tab:architecture}
\end{table}

Four levels of specification have distinct roles. The \emph{protocol} fixes
these objects, control boundaries and acceptance obligations. A \emph{profile}
chooses their exact encoding, proof language, checker procedure and permitted
foundation. The \emph{reference implementation} executes one profile. A
\emph{deployment} supplies the trusted checker installation, protected local
state and stable storage assumed by that execution. Changing a hash scheme
or proof representation requires a declared profile with its own acceptance
rules. It preserves the architecture when exact selection, independent proof
admission and the consumer's control remain intact. Unsupported profiles are rejected.

The adversary may supply a whole bundle, valid signatures, false interface
claims and arbitrarily nested publications. It may also alter received files
or present a forged acceptance record. The consumer's lock, checking software,
local authentication secret and accepted storage remain outside that authority.
This allocation is the basis for the acceptance procedure below. Compromise of
a consumer-controlled facility defeats its associated guarantee and requires
additional deployment controls.

\subsection{Assembly and immutable publication}

The provider chooses the theorem names to export, admits its own dependencies,
and builds the supplied source in a trusted publishing environment. It then
checks the compiled proof data and inspects the actual theorem interfaces and
axiom dependencies. Only a successful check enables the reference publishing
route to issue eligible declaration credits. Possession of a signing key is
unrestricted; signatures produced outside this route still face the full
consumer procedure.

The publication records a file inventory, a module-to-file map and the exact
identities of inherited publications. Its dependency binding includes both a
bundle identity and the certificate identities exported by that dependency.
Each declaration certificate binds its own interface and axiom policy to the
bundle. These relationships let a consumer distinguish a changed proof
realization, a different exported conclusion and an additional provider credit.
The source, proof and interface bindings change mathematical identity; another
valid credit preserves that identity. Mirrors may change download locations
without changing the consumer's selection.

\subsection{Consumer admission algorithm}
\label{sec:admission}

The inputs are a received bundle, an independently selected consumer lock, a
supported checking profile and consumer-controlled local state. Admission
processes the complete finite publication closure, meaning all publications
reachable through its dependency bindings. Shared components may be checked
once by identity. Every declared certificate remains covered, including an
inner conclusion that the chosen outer theorem does not use.

A conforming admission procedure performs the following checks in order:
\begin{enumerate}
\item \textbf{Select and validate the object.} Enforce the supported schema and
profile; recompute bundle and certificate identities; require the selected
bundle and certificates to match the consumer lock. Reject malformed objects,
unsupported versions and absent selections.
\item \textbf{Validate the whole publication closure.} Check all provider
signatures, every source and proof file, the module inventory, and each
parent-to-dependency binding. Inherited module bytes must agree wherever they
occur. Reject missing, unlisted or conflicting data before proof replay.
\item \textbf{Resolve local acceptance reuse.} If a local record exists,
authenticate it with the consumer's own secret. An invalid record fails.
Reuse requires a valid local record whose full checking scope matches the
current object and checking installation. A matching scope returns the reused
decision; a stale scope continues with the fresh-check steps. Incoming provider records cannot authorize this step.
\item \textbf{Check actual proof data.} Launch the trusted checker without
provider-controlled search paths or plugins. Read the complete proof data,
replay its declarations through the kernel, and inspect the actual theorem
kind, type, universe parameters and reachable assumptions for every certificate.
Compare the loaded module closure with the declared inventory.
\item \textbf{Compare each certificate separately.} Require its checked
interface to equal the claimed interface and its actual axiom set to lie
within that certificate's policy. Check its original proof dependencies against
its own realization; modules supplied only by an outer publication cannot fill
an omission in an inner one. Agreement with a broader policy used for
another certificate is insufficient. Every accepted foundational axiom must
have the identity required by the trusted foundation.
\item \textbf{Seal the decision for use.} Recheck the publication bindings and
input bytes after a fresh check. Record the checker result together with its
scope in authenticated local state. Make the derived build depend on acceptance
of these inputs; a failure returns an error and prevents a successful result
for the attempted build.
\end{enumerate}

The outputs distinguish \emph{content and credit verified}, \emph{fresh
mathematical acceptance}, and \emph{local acceptance reused}. The first result
covers object integrity and attribution. The latter two additionally cover an
actual proof check under the recorded scope, performed now or retained locally.
This distinction must remain visible to both the build and the reader of its
result. Both fresh and reused decisions remain subject to the same controlled-use
and build-dependency requirements. No transition constructs a theorem from a signature.

\subsection{Keeping acceptance valid through use}

A local acceptance record contains the checked interfaces and their policies,
source and proof inventories, the checking code and executable identities, the
loaded foundation-file identities, and the checker report. Provider credits
are verified on each admission request, but are outside the reusable
mathematical scope. Consequently, an additional valid credit needs no new proof
replay. A changed proof, interface, foundation or checker invalidates the
matching record and requires the appropriate new selection or fresh check.

The accepted files must remain stable between checking and use. The reference
runner acquires an exclusive admission lease and rechecks bytes before sealing
success. The lease coordinates cooperating local runs; storage controlled by
an attacker requires a read-only snapshot or equivalent protection. Similarly,
a digest records which executable was checked but does not prevent that
executable from being replaced afterwards. The deployment preserves the
installation and accepted store for the lifetime of the reused acceptance.

The normal build invokes this admission path even when its own source is
unchanged. It then uses the admitted declarations in an ordinary proof. A
successful derived publication binds the inherited bundle and certificate
identities again, so the next consumer can execute the same contract.

\vspace{0pt}\Needspace{20\baselineskip}
\section{Composition and the security boundary}

The composition argument separates logical validity from object binding.

\begin{proposition}[Composition]
% formal-claim-begin: composition
Let
$C$ be a type of claims, let $F,K,M:C\to\mathrm{Prop}$ denote foundation
membership, local checked inference, and intended validity, respectively, and
let $D:C\to\mathrm{List}(C)$ give the dependencies of a claim. Define
$\Accept(c)$ inductively: it holds when $F(c)$ holds, or when $K(c)$ holds and
$\Accept(d)$ holds for every $d\in D(c)$. This describes a finite acceptance
derivation, with repeated dependencies permitted.

For all such $C,F,K,M,D$, assume
$\forall c,\ F(c)\Rightarrow M(c)$ and
$\forall c,\ K(c)\Rightarrow
((\forall d\in D(c),\ M(d))\Rightarrow M(c))$.
For every root $r\in C$, $\Accept(r)$ implies $M(r)$.
% formal-claim-end: composition
\end{proposition}

\begin{proof}
% relation-claim-begin: composition-proof
Induct on the acceptance derivation. In the foundation case, the first
assumption gives validity. In the derived case, the induction hypotheses give
validity for every listed dependency; the second assumption combines these
facts with the local checked inference to give validity of the root.
% relation-claim-end: composition-proof
\end{proof}

There is no fixed depth in this argument. Its local soundness premise must
include faithful interface binding and complete dependency coverage. The
proposition does not discover missing dependencies or prove the correctness of
a parser, kernel implementation, signature library, or filesystem. The Lean
companion supplied with the reference implementation (Section~\ref{sec:lean})
states this same abstract induction theorem explicitly; applying it
to an implementation retains those implementation assumptions.

\begin{proposition}[Credit independence]
% formal-claim-begin: credit-independence
For any types $C,Q$, predicate $A:C\to\mathrm{Prop}$, claim $c\in C$, and
credits $a,b\in Q$, define mathematical acceptance by
$\mathsf{MathematicalAcceptance}(A,c,q):=A(c)$.
Then mathematical acceptance with credit $a$ is equivalent to mathematical
acceptance with credit $b$.
% formal-claim-end: credit-independence
\end{proposition}

\begin{proof}
% relation-claim-begin: credit-proof
Both sides are the same proposition $A(c)$ by definition.
% relation-claim-end: credit-proof
\end{proof}

\begin{proposition}[Replacement of valid credit]
% formal-claim-begin: valid-credit
For any types $C,Q$, predicates $A:C\to\mathrm{Prop}$ and
$V:Q\to\mathrm{Prop}$, claim $c\in C$, and credits $a,b\in Q$,
if $V(a)$ and $V(b)$, then
$(A(c)\land V(a))\leftrightarrow(A(c)\land V(b))$.
% formal-claim-end: valid-credit
\end{proposition}

\begin{proof}
% relation-claim-begin: valid-credit-proof
In either direction retain $A(c)$ and pair it with the assumed validity of the
replacement credit.
% relation-claim-end: valid-credit-proof
\end{proof}

The latter two propositions record a deliberate separation of roles; their
simple proofs reflect that design. Protocol-level eligibility may require a
valid credit, while mathematical validity remains independent of who provides
it. A compromised signing key can misattribute supply. Under an unchanged
consumer lock it cannot silently replace the bound proof, and it cannot make a
false type pass a sound proof checker even if a consumer approves a new lock.

Replacement detection assumes collision resistance of the selected digest and
correct verification of signatures. A changed proof, source, interface, or
dependency changes the bound identity. Re-signing a changed graph does not
restore its former identity. A genuinely valid replacement may be accepted
under a new lock, preserving integrity through an explicit new selection.
The registry may withhold objects or present conflicting discovery views; these
are availability and discovery problems. The accepted identities remain
consumer-selected.

\section{Concrete Lean profile and assurance argument}
\label{sec:lean}

The reference profile fixes Lean 4.30.0-rc2 and supplies an export/admission
library, command-line operations and a default Lake target. Lake is Lean's
project build tool; the default target is the operation run by an ordinary
\texttt{lake build}. The code and normative specification are located in the
\hyperref[sec:artifacts]{artifact availability statement}. This section gives
the technical construction that implements Section~\ref{sec:admission}.

\subsection{Encoding and cross-object commitments}

The \emph{core} of an object is the set of fields committed by its identity.
The bundle core contains the profile and toolchain, module map, file hashes
and dependency bindings. A declaration core contains its bundle identity,
interface and allowed axioms. The full transport envelope also carries credits
and nested publications, which are checked against those cores.

Let $\operatorname{can}(x)$ be the canonical byte encoding of core $x$.
This profile uses UTF-8 JSON with no inter-token whitespace, preserves literal Unicode
scalar strings, orders object keys by UTF-8 bytes and preserves array order.
It permits null, booleans, strings, arrays, objects and integers in
$[-(2^{53}-1),2^{53}-1]$. Duplicate keys, floating-point values and invalid
Unicode are rejected. This gives independent readers the same byte sequence
without silently normalizing a name or expression.

SHA-256 supplies the content hash. For an object kind labelled $d$, its identity is
\[
\operatorname{id}_{d}(x)=\texttt{sha256:}\,\Vert\,
\operatorname{hex}\!\left(\operatorname{SHA256}
  (\operatorname{tag}_{d}\Vert\texttt{0x00}\Vert\operatorname{can}(x))\right).
\]
Here $\Vert$ denotes byte concatenation, \texttt{0x00} is one zero byte,
$\operatorname{hex}$ renders a digest
as hexadecimal, and $\operatorname{tag}_{d}$ is the UTF-8 string
\texttt{theorem-notary/}$d$\texttt{/v1}. Distinct domains for bundles and
components prevent one kind of object from being reinterpreted as the other.
A provider signs a separate credit body containing the component identity,
provider role, declared name, public key and key fingerprint. The Ed25519
signature scheme signs that canonical body with the separate
\texttt{theorem-notary/provider-credit/v1} tag and a zero-byte separator.
The public key identifies the cryptographic supplier; a human-readable name
remains provider-declared.

The interface fields are the declaration's full name, structural type
expression and universe parameters. The pinned Lean expression representation
retains referenced constants and binder information; the bundle fixes the
definitions that give those constants meaning. Source and proof paths are
relative to the bundle, and each has an exact digest. The profile binds the
compiled module data, including exported, server and private proof-data parts
when present, so a hidden part cannot supply an unlisted dependency.
Compiled realizations are version-specific; profile migration requires an
explicitly justified new realization.

\subsection{Data-only checking and complete dependency coverage}

The checker starts as a separate process with trusted search paths. It first
loads its own checking code and the official foundation. Only then does it
introduce the supplied proof-data directory. Before launch, the runner rejects
supplied modules that would shadow installed toolchain modules. The checker
disables imported extensions and provider plugins, reads proof data, and replays
declarations into a fresh kernel environment. Thus the provider's code cannot
obtain a successful admission merely by printing the expected output or
running a module initializer.

To preserve each publication's original proof realization, the checker also
compares the declaration records retained in the individual imported modules.
Where a supplied module shares a declaration name with another module, the
records must agree structurally, including the proof body. This check precedes
reliance on the merged environment: a same-type theorem with a different proof
cannot discharge an earlier component's unproved obligation. Each inherited
module must also retain exactly the same ordered proof-data parts in its parent.

For each certificate, the object must be a theorem in the original proof data.
Its type and universe parameters must agree with the replayed declaration.
The checker traverses the constants actually used in its type and proof,
recursively inspects their definitions, and records the reachable axioms and
referenced theorems. An uncovered constant fails. The runner also compares the
actually loaded non-foundation modules with the complete bundle inventory;
module parts must come entirely from either the bundle or the trusted
installation. Uncontrolled paths, missing modules and conflicting data fail.

Coverage is also checked within each inherited realization. The checker retains
the original module origins of reachable constants, including dependencies
through field projections, mutually defined inductive types and their
recursors, the rules for eliminating inductive values. For each
such constant, at least one identical original occurrence must belong to that
certificate's realization or the trusted installation. Multiple identical
origins are alternatives, so import order does not assign arbitrary ownership.
Each module's recorded imports must likewise be available within that
realization. The outer union of valid modules therefore cannot conceal an
inner publication's missing dependency. These comparisons allow one replay to
support several certificates while preserving each certificate's own scope.

The permitted foundational axioms are \texttt{propext}, which supplies
propositional extensionality; \texttt{Classical.choice}, which supplies classical
choice; and \texttt{Quot.sound}, which supplies quotient soundness. Their actual
types and universe parameters must match the official installed foundation.
A certificate may allow a smaller subset. Even when one replay covers several
certificates under their combined axiom set, the runner separately enforces
each certificate's narrower policy. Reachable unproved assertions introduced by \texttt{sorry} or extra
project axioms, an axiom offered as a theorem, and unsafe or partial reachable
proof dependencies fail. The composition companion and the small reuse case
require none of these foundational axioms.

\subsection{Local records and build behavior}

Fresh admission produces a structured checker report listing the replay state,
extension-initialization state, loaded module parts, checked interfaces,
actual axioms and referenced theorems. The runner records this report with the
scope described in Section~\ref{sec:architecture}. It authenticates the local
record using HMAC-SHA-256, a keyed message authentication code whose secret is
held by the consumer outside the publication. This local format is a cache
control for the reference runner; public certificate exchange does not require
sharing that secret.

The admission operation returns \texttt{created} for fresh acceptance or
\texttt{reused} for a matching local record, together with the bundle identity,
replay count and record location. Inspection alone reports a content-and-credit
assurance label. The generated default Lake target invokes the build runner,
which admits dependencies before producing a derived result. It revisits those
dependencies even when reusing an unchanged derived publication. A failed
invocation returns a nonzero exit code and replaces the attempted build's
result with a failure state. A stale successful result therefore cannot mask
the rejection in this tested build path.

\subsection{Why these mechanisms establish the claimed assurance}

The assurance argument follows the same object through successive checks.
Canonical commitments and the independently held lock fix which publication
is being evaluated. Complete inventories prevent extra proof files from
entering unnoticed. Kernel replay establishes the formal content of the proof
data; structural comparison ties that content to each advertised declaration.
Original declaration comparison preserves the proof objects across import merging.
Original-origin and import checks preserve each realization's dependency scope.
Per-certificate axiom checks enforce the selected foundation. Local record
authentication preserves the scope of a past check, and storage control
preserves its inputs through use. Table~\ref{tab:assurance} gives the decisive
check for representative failures. These are separate, falsifiable obligations;
passing one does not grant the guarantees of the others.

\begin{table}[ht]
\centering\small
\renewcommand{\arraystretch}{1.12}
\begin{tabularx}{\linewidth}{@{}>{\raggedright\arraybackslash}X>{\raggedright\arraybackslash}X@{}}
\toprule
Failure presented to the receiver & Decisive check and supporting assumption\\
\midrule
Changed proof or dependency, consistently re-signed & Recomputed identities must match the consumer's existing lock; relies on digest collision resistance and protected selection.\\[3pt]
False interface under an approved new lock & Actual replayed type must match every certificate; relies on checker soundness and complete dependency inspection.\\[3pt]
A valid outer theorem carrying a false inner certificate & All inherited certificates receive interface and axiom comparison, including unused exports.\\[3pt]
Same theorem name carrying a different proof body & Original per-module declaration comparison and identical inherited part inventories preserve the checked realization.\\[3pt]
Inner dependency omitted but supplied by an outer publication & Original constant origins and actual imports must be covered by the inner realization or trusted foundation.\\[3pt]
Provider code impersonating the checker & Trusted startup paths, module-shadow rejection and disabled import extensions isolate the tested admission path.\\[3pt]
Forged record or changed inputs after acceptance & Consumer-owned authentication and exact scope comparison reject forged or stale reuse; stable storage must hold through use.\\
\bottomrule
\end{tabularx}
\caption{Technical assurance is attached to explicit checks and their trusted
facilities. The adversarial evaluation exercises these boundaries.}
\label{tab:assurance}
\end{table}

This procedure is independent of the provider's private services: a fresh
consumer with the public bundle, its selected lock and the trusted checking
installation can perform admission. The first profile uses Lean's existing
binary reader and kernel. The abstract composition theorem assumes the
soundness and faithful binding of these local checks; it does not verify the
runner, cryptographic libraries or binary reader. Resource control limits
nested publication envelopes to 64 levels. Public conformance vectors pair specified inputs with expected encoding,
identity and signature results. Python and Node.js produce the same committed
bytes and verify the same signatures for those inputs.

Source correspondence and deployment protection have their own controls.
The publishing route compiles the supplied source, while the receiving checker
validates the resulting proof data. Independently rebuilding the readable
source would additionally test that it produces those data. The present
profile assumes a trusted publishing build environment and an immutable
checking installation, including its runtime dependencies. The runner directly
hashes its own code, checker source, Lean executable and loaded foundation
parts; those hashes do not inventory every executable dependency on the host.
Production acceptance of hostile source uploads requires an execution sandbox,
and hostile concurrent file writers require protected snapshots. The profile
makes these responsibilities visible at the deployment boundary.

\section{Publication and reuse cases}

\subsection{An independently published mathematical component}

The large-component case has a narrow theorem interface backed by 30,638
project modules and about 121 GiB of unpacked proof data. It uses an archived
formalization of the exact extremal bound in Erdős Problem 848. The result
concerns the maximum size of
$A\subseteq\{1,\ldots,N\}$ for which $ab+1$ is nonsquarefree for every
$a,b\in A$, including $a=b$; its exact bound for positive $N$ is
\mbox{$\lfloor(N+18)/25\rfloor$}~\citep{erdos848}. The protocol experiment uses this existing formal result as a publication
case; its contribution concerns the handoff of the proved statement.

The published certificate binds the actual all-$N$ declaration, its fixed
source release, the proof archive, and the allowed foundation. Its provider
credit names Alex Chengyu Li and a public signing key. The name is
provider-declared; signature verification establishes key attribution, not
independent verification of the person's identity or priority. An external
consumer imports the original theorem and proves its own literal alias, while
exposing the certificate's provenance and checking basis. A consumer choosing
this result is selecting the full positive-integer range, the diagonal pairs,
and the exact extremal bound. Reading the mathematical argument establishes
why those features matter. The certificate binds that selection to the actual
declaration and proof archive, so later consumers can identify the supplied
formal result and inspect its checking history.

Restoring this case required about 30 GiB of archive downloads. Its historical
kernel archive is accepted under an explicit, immutable archive policy, and
the new external consumer was checked separately. The experiment does not
claim a fresh complete replay of the upstream proof under the new generic
profile. These two assurance histories remain visible instead of being
collapsed into one label.

\subsection{Derived publications across three projects}

A separate fresh-proof case tests composition without relying on the large
historical archive. A provider defines binary trees and proves two conclusions:
mirroring twice returns the original tree, and mirroring preserves leaf count.
Each receives a declaration certificate. A second project imports both,
packages double mirroring as a round trip, and proves that the round trip
preserves the tree and its leaf count. A third project mirrors the round trip
once more and defines a frame size as its leaf count plus a fixed overhead of
four. Using both middle-project conclusions, it proves that the output equals
the original tree's mirror and that the frame size equals the original leaf
count plus four. These are six distinct
certificates carried by three projects. The second project is both a consumer
and a new provider. The third project receives its conclusions together with
the inherited publication identities, and can check that whole chain without
taking the second provider's account of earlier checks as sufficient evidence.
This is the protocol's recurring use case: a researcher publishes a result
that remains usable when its consumer subsequently becomes a provider.

The final proof dependency reports identify both of the middle project's
theorems. Thus the example exercises actual theorem use, not only the presence
of signed references beside an unrelated proof. The middle project's default
Lake target performs admission before producing its result. Modifying the
imported proof causes this ordinary build to fail. A separate eight-layer
fixture, with two conclusions at every layer, exercises deeper wrapping; the
abstract composition argument is independent of those tested depths.

The receiving context is a separate process supplied with public bundle bytes,
a consumer-selected lock, fresh acceptance state, and a different local key.
Its first admission replays the proof. Repetition and the addition of a second
valid provider credit reuse the mathematical acceptance record. A source
revision, even one preserving the theorem statements, receives a new identity;
the old lock rejects it, and an explicitly updated lock triggers new admission.
This automated cold context is on the same physical workstation. Reproduction
by another operator or on another platform remains a further validation step.

\section{Adversarial evaluation and cost}

The evaluation asks whether a receiver can be made to accept a different
statement, lose an inner component's proof obligation, or reuse a fabricated
checking result. The adversary has the bundle and file controls of
Section~\ref{sec:architecture}, including its own valid signing keys. The
consumer controls its lock, local authentication secret and checking installation.
The cases below specify the attempted deception and the observed rejection.
The executable inputs and complete result inventory are located by
\texttt{Notary/REFERENCE.md} and \texttt{Notary/reference-result.json} in the
public repository listed under Artifact availability.

\subsection{Replacing the selected theorem or its inner justification}

A valid signature is particularly revealing when it authenticates the wrong
claim. We replace a tree theorem's advertised type with \texttt{False}, leave
its compiled proof unchanged, and sign the modified certificate with a valid
provider key. The original consumer lock rejects the changed certificate
identity. We then explicitly select the altered certificate in a new lock,
so that this identity check no longer decides the outcome. Admission still
rejects it: the advertised \texttt{False} type differs from the declaration
actually checked by the kernel. Placing the same false certificate inside a
derived publication, recomputing the outer dependency bindings, and signing
the outer certificates also ends in this interface rejection. The checks
therefore cover the inherited claim even when the outer proof remains valid.

The original proof presents a different substitution opportunity. Two compiled
modules export the same named theorem of \texttt{True}; one uses
\texttt{sorry} to leave its proof obligation unresolved, while the other uses
\texttt{True.intro}. A merged
environment could present the valid proof for that name and thereby conceal
the first component's unresolved obligation. The receiver rejects the pair
when it compares the original module declarations, before accepting that
merged view. Checking the module containing \texttt{sorry} alone also fails for its
\texttt{sorryAx} dependency. The proposition is true in this example; the
rejection preserves which component actually supplied a complete proof.

To challenge dependency ownership, we remove the tree provider from the inner
round-trip publication's dependencies and module inventory. We retain the
provider's proof files in the outer application and re-sign the changed inner
and outer publications. Content and signature checks accept the internally
consistent descriptions. Fresh mathematical admission rejects the inner
round-trip module's actual import, which its own realization no longer covers.
The outer application's possession of a valid proof therefore does not repair
the inner publication's missing dependency. This case locates the required
check at each realization's boundary, beyond verification of the outer union.

A permitted update supplies a comparison with these failures. Appending a source
comment preserves the theorem statements but changes the source realization.
The old lock rejects it; an explicitly selected new lock leads to fresh
admission. Adding another valid provider credit to an unchanged component
instead reuses the mathematical acceptance record. Together these observations
distinguish a new component choice from a change in who supplies it.

\subsection{Protecting the checking process and the build decision}

We flip a byte in an imported compiled module and attempt to reuse the earlier
acceptance. Inventory validation rejects the file before kernel replay. Running
the consumer's ordinary \texttt{lake build} against that dependency also fails
and replaces the stale successful build result with a failure state. Adding a
compiled module absent from the inventory is rejected as unlisted data. These
observations connect the public binding rule to the operation a downstream
author actually invokes.

To impersonate trusted checking code, we place provider-controlled files at the
paths of \texttt{Lean.Replay} and \texttt{Init}; the runner rejects both before
launching the checker. A separate module has a valid theorem together with an
initializer that writes a marker file. Ordinary import writes the marker,
providing a control that the code can execute. Admission of the same proof
data leaves the marker absent. Another source prints a success-looking message
but offers a proof of \texttt{False} through \texttt{sorry}; the independent
checker rejects its axiom dependency and the publisher produces no certificate.
These experiments exercise search-path isolation, disabled initialization and
actual proof inspection. The binary reader's memory safety remains a trusted
implementation property.

Finally, we edit the stored checker report inside an existing local acceptance
record while retaining its old authentication tag. Reuse fails authentication.
Supplying an empty or malformed consumer lock also fails, even with a valid
previous record available. Thus possession of cached success cannot replace
selection, and changing a cached report requires the consumer's authentication
secret. Lower-level regression cases check that duplicate JSON keys and invalid
numeric encodings are rejected, distinct object domains have distinct
commitments, credit identity stays bound to its key, and inherited proof-data
parts preserve their exact ordering and membership. Their purpose is to preserve
the specified input contract between revisions; the public inventory identifies
the particular behaviors exercised.

\subsection{Cost of fresh acceptance and subsequent reuse}

The cost question is how much checking remains when a consumer receives a
component for the first time, uses it again, or explicitly adopts changed source.
Table~\ref{tab:cost} records these operations alongside content inspection and
ordinary builds, whose assurance scopes differ.

\begin{table}[ht]
\centering\small
\begin{tabularx}{\linewidth}{@{}Xr@{}}
\toprule
Operation & Elapsed seconds\\
\midrule
Fresh consumer mathematical admission & 64.8\\
Unchanged consumer acceptance reuse & 3.0\\
Updated dependency with an explicitly new lock & 69.5\\
Content and credit verification alone & 0.005\\
Plain Lake build with fresh project state & 1.8\\
Plain Lake warm build & 0.9\\
Notary-consuming Lake warm build & 7.1\\
\bottomrule
\end{tabularx}
\caption{Single observations from the small reuse experiment. The operations
have different assurance scopes. An equal-assurance comparison of performance
with ordinary cached imports remains a separate evaluation task.}
\label{tab:cost}
\end{table}

Table~\ref{tab:cost} distinguishes costs that a single successful-build label
would hide. Fresh admission includes whole-environment replay. Reuse verifies
the acceptance scope and current bytes. Plain Lake measurements provide an
ordinary build baseline, whereas content-only checking supplies no new
mathematical evidence. The measurements were made on a Windows workstation
under normal load without resetting the operating-system file cache. Cold
means fresh project or acceptance state. The operating-system cache remained
available, and network downloads were outside the timed operations. Successful publisher and derived-package artifacts
were retained during the completed evaluation; their reuse is recorded
separately from fresh consumer admission.

The result is useful acceptance reuse, accompanied by substantial initial
checking overhead. It motivates research on checked persistent environments,
proof-closure deduplication, and declaration-level reuse under equivalent
assurance. Those optimizations must preserve exact bindings and coverage. A
smaller downloadable certificate or a cryptographic signature cannot by itself
replace the proof obligation.

\section{Relation to existing systems and implications}

OpenTheory provides directly relevant precedent for packaging logical
theories, tracking dependencies, and exporting and importing them across higher-order logic (HOL)
systems~\citep{hurd2009}. The present Lean profile has narrower logical
interoperability. Its focus is an explicit contract joining declaration
selection, provider credit, consumer-controlled admission, and derived
publication. Existing proof exchange is the starting point for this
specialization.

Software supply-chain systems provide a second precedent. In-toto binds
artifacts and development stages cryptographically so that consumers can
check the integrity of a software supply chain~\citep{intoto2019}. The theorem
setting requires an additional distinction: an authentic process record or
provider credit is insufficient for mathematical admission. The actual
declaration must pass the relevant logical checker under its disclosed
foundation. Conversely, a correct theorem does not by itself authenticate
provider identity, prove source correspondence, or ensure registry
availability.

The practical gain is continuity between mathematical choice and subsequent
use. A reader can first decide whether a theorem says what is needed, then
select its exact publication, and finally build a new result whose inherited
basis remains inspectable. A collaborator receiving that result obtains
explicit objects and admission conditions in place of an implicit agreement
with every earlier maintainer. The protocol thereby makes careful theorem
provision a reproducible role in a larger research collaboration.

Adoption can begin at one publication boundary. A project may keep its ordinary
Lean proofs and use the admission target when consuming a published component.
Content-addressed publications can be distributed through files, repositories,
or mirrors; a central registry adds discovery and distribution. The provider's
credit remains visible across that distribution, and consumers continue to
choose their own locks and checking conditions. This organization accommodates
independent providers alongside shared libraries, with publication granularity
at the level of each eligible theorem.

The strongest reason to adopt the protocol is a need to carry this explicit
contract across maintainers, releases, or derived publications. An existing
workflow that already supplies the same bindings and checks can implement the
contract with familiar tools. The reference implementation demonstrates that
the contract can be executed and tested. Independent consumers, further
substantial components, and equal-assurance workflow comparisons can establish
where standardization saves effort in practice. The current timings expose
the cost of fresh admission and the benefit of reusing its unchanged local
result; theorem interpretation, source correspondence and version migration
retain their own obligations.

The long-term objective of a comprehensively formalized mathematical library
makes this contract especially consequential. Such a library would give
subsequent research a growing body of explicitly stated, checkable inputs.
Automated iteration requires a machine to distinguish an accepted component,
a changed realization, and a candidate still awaiting proof. Per-declaration
publication and explicit admission make those distinctions available to a
build process. New mathematical work can then accumulate on a disclosed basis,
with formalization effort directed toward extending that basis. The reference
profile establishes one part of this organization; constructing and reviewing
the new proofs remains substantive work.

A theorem publication centre built on this protocol would therefore organize
access to independently checkable mathematical components. Providers receive
credit for making those components available. Researchers preserve control over
which theorem and foundation enter their work. Derived publications retain
that discipline at every handoff, allowing the output of one proof project to
become an explicit, inspectable input to the next.

\section*{Artifact availability}
\label{sec:artifacts}
The reference implementation, normative profile, public wire vectors, and
reproduction instructions are available at
\url{https://github.com/crabsatellite/mathlib-theorem-notary}.
The repository's \texttt{Notary/reference-result.json} identifies the evaluated
source files by digest and records the experimental results. The large \#848
source and proof archive are available from
\url{https://github.com/crabsatellite/erdos-848-squarefree-product}.
Local acceptance secrets are not part of either public artifact.

\section*{AI assistance}
OpenAI Codex was used as an AI-agent component in software development and
manuscript drafting.
The author is responsible for the protocol claims and their interpretation;
the reported checks refer to the recorded artifact executions.

\bibliographystyle{plainnat}
\bibliography{paper_references}
\end{document}
